\paragraph{From the incidence origin to lattice modification.}
\label{inc-ampl-reticulo}
The marked flag determines a position through an equivariant
set operation:
\[
 o_*=(H_e\cap H_\varphi)
       \setminus(H^-_\alpha\cup H_{\rm APP})=\{1\}.
\]
Its function is to distinguish one component among the twelve entering
the lattice realization. The complete ternary code
acts as a glue code in the discriminant module
$(A_2^*/A_2)^{12}$. Its preimage defines $N(C_W)$, and
self-duality, weights, and distance of the code translate,
respectively, into properties of integrality and determinant, parity,
and norm bounds for the glue classes. Here the symbolic information
of the code performs a function no longer performed by the isolated support.

The $A_2$ metric and positive radial chamber are fixed in the lattice
development. Within that chamber, the congruence required for an even
index-three neighbor has twelve realizations of minimum norm $54$,
permuted by changing the distinguished component. The origin $o_*$
selects
\[
 v_\alpha=4\rho_1+\rho_2+\cdots+\rho_{12},\qquad
 v_\alpha^2=2(4^2+11)=54.
\]
The coefficient $4$ belongs to that solution of the congruence in the
declared chamber. The subscript $\alpha$ records the incidence
provenance of the selection: its content is the frame that entered
closure and now individualizes a lattice direction.

Pairing with $v_\alpha$ determines the index-three sublattice
$N_{\alpha,0}$, and adding the class
$v_\alpha/3$ produces the neighbor $N_\alpha$ of
\eqref{exc:vecino}. This modification preserves rank and recovers
unimodularity and parity; congruence selection eliminates the previous
roots, and local bounds exclude roots in the new classes.
Identification with Leech appears after these verifications. The
binary realization of the five regions and this lattice construction
are two prolongations of the common frame: the second receives the
ternary code directly as glue data.
